\scanpage{71}
plus intrinsèque, le graphe $G_0$ \deleted{\unsure{se déploie}} \unsure{associe|associé} à l'ens $I\amalg\widetilde A'_0$ de card $4$, muni d'une partition de type $(2,2)$, (et \unsure{$=$} d'une partie \inserted{$I$} (de card $2$), l'ens \inserted{$J$} des \inserted{$3$} partitions de type $(2,2)$ a donc un pt marqué $\varepsilon_1$, les deux autres pts correspondent aux élts de l'ens $J':=I\wedge_{\mathbb Z/2}\widetilde A'_0$ :
\[
(84)\qquad J=\{\varepsilon_1\}\amalg J'\qquad J'=I\wedge_{\mathbb Z/2}\widetilde A'_0.
\]
Ceci posé, $\mathbb P^1_J=\mathbb P^1_{\{\varepsilon_1\}\amalg J'}\deleted{\unsure{\mathfrak S_{J'}}}$ peut être considéré comme le compactifié de $\mathbb G_m^{\unsure{1}}\wedge_{\pm1}J'$ et \deleted{$\unsure{\mathfrak S_2}$} $\mathfrak S_{J'}\simeq\mathbb Z/2$ y opère par transfert de structure. On a :
\margin{$\varepsilon_1$ correspond au point $1$}
\[
(85)\qquad
\left\{\begin{aligned}
\widehat M_{G_0}&\simeq(\mathbb P^1_{\{\varepsilon_1\}\amalg J'}\deleted{\unsure{\mathfrak S_{J'}}})\\
(\widehat M_{G_0},\unsure{\Gamma'_0})&\simeq(\mathbb P^1_{(\varepsilon_1\amalg J')},\mathfrak S_{J'})
\end{aligned}\right.
\]
\margin{$(\unsure{\widehat M_{D_0\bar\varepsilon}})$}
sur lesquels $\mathfrak S_I$ opère via $\mathfrak S_{J'}$, donc $\mathfrak S_I$ \uline{opère trivialt sur} $\widehat M_{D_0}$. Les deux pts à l'infini de $\widehat M_{D_0}$ correspondent : la trajectoire $J'$ d'une part \inserted{(pt. $P_0$)}, la trajectoire $\varepsilon_1$ d'autre part (point $P_1$). On a
\[
(86)\qquad M_{D_0}\simeq\unsure{\mathbb E^1_{\mathbb Z}|\mathbb A^1_{\mathbb Z}}\setminus\{\text{section }2\}=\operatorname{Spec}\mathbb Z[T]_{(T-2)}
\]
\[
\text{via}\quad\left\{\begin{aligned}M_{G_0}&\longrightarrow M_{D_0}\\ T&\longrightarrow T+\frac1T\end{aligned}\right.
\]
\[
(87)\qquad\widehat M_{D_0}\simeq\bigl(\mathbb P^1_{\mathbb Z},\text{ avec ``signature'' }2\text{ le long de la section }T=2\bigr)
\]
\margin{70}
